% =============================================================================
% Section 3 -- A Taxonomy of Agent Failure Modes
% Contains Table 1 (tab:taxonomy), a single-column table.
% =============================================================================
\section{A Taxonomy of Agent Failure Modes}
\label{sec:taxonomy}

A critical prerequisite to structured remediation is structured cause-finding.
We propose a five-category taxonomy of agent failure modes relevant to
customer-facing agents: four intrinsic modes---behavioral drift, input
brittleness, alignment creep, and tool misuse---plus one exogenous mode,
mesh disruption (\texttt{exogenous\_disruption}), in which the failing agent is
not the cause of its own degradation
(Section~\ref{subsec:mesh}). These categories are not mutually exclusive, but they
suggest distinct remediation pathways, which is the operationally important
property for APIP design. Following Kirkpatrick's~\cite{kirkpatrick1994evaluating}
insight that evaluation should be hierarchical, we note that each failure mode
category corresponds to a different level of the agent's behavioral stack, and
consequently to a different intervention tier.

\begin{table}[t]
  \centering
  \caption{Agent failure mode taxonomy with Kirkpatrick level mapping and
           remediation tier.}
  \label{tab:taxonomy}
  \footnotesize
  \begin{tabular}{@{}L{1.4cm}L{3.6cm}L{2.1cm}@{}}
    \toprule
    \textbf{Failure Mode} & \textbf{Description \& Signals} &
    \textbf{Remediation Tier} \\
    \midrule
    Behavioral Drift &
    Gradual degradation as world-state diverges from training or retrieval
    context. Signals: rising hallucination rate on new entities, outdated
    policy citations, product terminology mismatches. Kirkpatrick Level 4
    (Results) failure. &
    Tier 2: Retrieval / RAG update \\
    \addlinespace
    Input Brittleness &
    Strong performance on modal inputs, catastrophic failure on a specific
    subclass. Signals: bimodal CSAT distribution, failure clustering around
    user segment or topic. Kirkpatrick Level 2 (Learning/Capability) failure. &
    Tier 1 (examples) or Tier 3 (fine-tune) \\
    \addlinespace
    Alignment Creep &
    Subtle drift in tone, value expression, or policy compliance. Signals:
    escalating human review flags, regulatory keyword triggers, brand
    guidelines violations. Kirkpatrick Level 3 (Behavior) failure. &
    Tier 1: Prompt constraint injection \\
    \addlinespace
    Tool Misuse &
    Correct intent identification with incorrect tool invocation. Signals:
    action errors, silent failures in downstream integrations, over- or
    under-use of tool surface. Kirkpatrick Level 3 (Behavior) failure. &
    Tier 1 (prompt) or Tier 3 (fine-tune) \\
    \addlinespace
    Exogenous Disruption &
    Peer agents degrade after a mesh change event (e.g.\ a new agent enters);
    no intrinsic failure mode found in the degraded agents. Signals:
    degradation concentrated in peer roles, onset coincident with a
    mesh-composition event. Kirkpatrick Level 4 (Results) failure at the
    system level. &
    Mesh protocol (Sec.~\ref{subsec:mesh}), not an agent-level tier \\
    \bottomrule
  \end{tabular}
\end{table}

\textbf{Behavioral Drift} is perhaps the most common failure mode in production
deployments. A customer service agent trained or prompted against a product
catalog from six months ago will increasingly produce incorrect citations as
the catalog evolves. This is not a model capability failure---it is a retrieval
or context freshness failure, and its remediation pathway is correspondingly
different from fine-tuning. In Kirkpatrick's terms, this is a Level 4 failure:
the agent's learned behavior has not changed, but the organizational results
are degrading because the environment has shifted.

\textbf{Input Brittleness} is particularly pernicious because aggregate metrics
mask it. An agent with 85\% task success overall may be failing catastrophically
on the 8\% of requests that involve a particular billing scenario or a
non-native-English speaker segment. Without disaggregated evaluation, the APIP
trigger will not fire until overall performance has substantially degraded.
This maps to a Kirkpatrick Level 2 failure: the agent lacks the capability
(learning) to handle a specific input class, even though its general capability
is adequate.

\textbf{Alignment Creep} covers the class of failures in which the agent's
behavior drifts from the intended value envelope---not necessarily producing
factually incorrect outputs, but producing outputs that are off-brand, subtly
non-compliant, or misaligned with the organization's customer communication
policies. Hendrycks et al.~\cite{hendrycks2021values} characterize this class
of failure as value misalignment, distinct from capability failure, and argue
it requires distinct detection and remediation mechanisms. In Kirkpatrick's
terms, this is a Level 3 failure: the agent's behavior has changed in ways that
do not serve the organizational goals, even though its underlying capability
has not degraded.

\textbf{Tool Misuse} is increasingly relevant as agentic systems are given
access to action surfaces: ticket creation, refund processing, escalation
routing, and external API calls. An agent that correctly identifies user intent
but invokes the wrong tool, invokes a tool with incorrect parameters, or fails
to invoke a necessary tool produces a qualitatively different class of failure
than a purely conversational agent---one with potentially significant downstream
consequences.

\textbf{Exogenous Disruption} is the taxonomy's escape hatch: the case in which
cause attribution over a degraded agent correctly concludes that \emph{no}
intrinsic failure mode explains the degradation, because the cause lies in a
change to the agent's mesh environment---most prominently the introduction of a
new, individually competent peer agent. Its signature is structural rather than
behavioral: degradation concentrated in peer roles, coincident with a
mesh-composition event, with the degraded agents' own traces annotating clean.
Because no agent-level intervention addresses the cause, its remediation is a
mesh-level protocol rather than a tier
(Section~\ref{subsec:mesh}); misclassifying it as an intrinsic mode and
applying an agent-level fix is precisely the tier-mismatch error the APIP is
designed to prevent.

\subsection{From Trace Labels to Taxonomy Categories}
\label{subsec:mast-mapping}

The taxonomy is only operational if a degradation window can be classified into
it by procedure rather than intuition. We publish the deterministic mapping we
use (Table~\ref{tab:mast-mapping}): cause attribution combines a trace-level
failure label distribution from MAST~\cite{cemri2025mast} with a
\emph{disaggregation signature}---how the degradation distributes across task
categories, agent roles, and time---to select an APIP failure mode and
intervention tier. The mapping is frozen before any evaluated run
(Section~\ref{sec:simulation}); below-confidence combinations \emph{abstain},
and abstentions are logged and scored as attribution failures rather than
silently resolved.

\begin{table}[t]
  \centering
  \caption{Published MAST-to-APIP mapping. Note the two rows sharing a MAST
           cluster but differing on disaggregation signature: trace labels
           alone cannot select a tier, which is why the taxonomy is
           two-level.}
  \label{tab:mast-mapping}
  \footnotesize
  \begin{tabular}{@{}L{2.2cm}L{2.6cm}L{1.5cm}L{0.7cm}@{}}
    \toprule
    \textbf{MAST cluster} & \textbf{Disaggregation signature} &
    \textbf{APIP mode} & \textbf{Tier} \\
    \midrule
    Verification / termination failures &
    Uniform across categories; tool-error rate elevated &
    \texttt{tool\_misuse} & 1, esc.\ 3 \\
    \addlinespace
    Specification / design failures &
    Concentrated in one task category or segment &
    \texttt{input\_brittleness} & 1, esc.\ 3 \\
    \addlinespace
    Specification / design failures &
    Uniform; entity/citation errors rising over time &
    \texttt{behavioral\_drift} & 2 \\
    \addlinespace
    Inter-agent misalignment &
    Concentrated in \emph{peer} roles after a mesh change event &
    \texttt{exogenous\_disruption} & mesh \\
    \addlinespace
    (no dominant MAST mode) &
    Slow uniform decline; review/flag rates falling &
    \texttt{alignment\_creep} & 1 \\
    \bottomrule
  \end{tabular}
\end{table}

A further category, \textbf{Adversarial Subversion}, captures failures whose
proximate cause is an adaptive attack rather than capability or value drift.
Recent work demonstrates that frontier models can embed prompt injections in
their outputs that consistently evade diverse monitors, and that
resampling-based protocols can amplify rather than suppress such
injections~\cite{schoen2025stress,bhatt2025ctrlz}. We treat adversarial
subversion as a separate, out-of-scope category because its remediation
pathway diverges from the five above: prompt revision, retrieval refresh, and
fine-tuning each address
symptoms but not the underlying adversarial dynamic. Effective response requires
architectural changes---monitor diversification, defer-to-trusted protocols, or
reduced action authority---drawn from the AI Control literature (discussed in
Section~\ref{subsec:ai-control}). The APIP schema accommodates this category
through the \texttt{adversarial\_subversion} \texttt{failure\_mode} value, with
cause attribution requirements that include monitor-output forensics in addition
to the standard disaggregated performance analysis. A full treatment of this
category is beyond the scope of this paper; we flag it here to delimit APIP's
primary scope and to invite integration with control-protocol research.
